\section{From two stable actions to a finite module}\label{conv:section}

The last step of the construction uses a finite bimodule $F$ and a stable
map $v:X\to F\otimes_A X$.  Two actions on maps out of $X$ will determine
all positive self-extensions of a module over a triangular algebra.  We
prove this conversion independently of the particular bimodule.  The
module itself is a finite quotient; complete resolutions are used to
compute its extensions.

\subsection{The conversion principle}

For a finite bimodule $F$ projective on both sides, tensoring with $F$
is exact and preserves projectives.  Over a symmetric algebra it
therefore induces an exact functor on the stable category and commutes
with shifts.  In the following statement we write $FX=F\otimes_A X$
and use these shift identifications when writing $F(g)$.

For the triangular algebra $\Lambda=\begin{pmatrix}A&0\\F&A\end{pmatrix}$,
we write a left module as a triple $(M,N,\eta)$: here $M,N$ are left
$A$-modules and $\eta:F\otimes_A M\to N$ is $A$-linear.  The action is
\[
 \begin{pmatrix}a&0\\f&b\end{pmatrix}(m,n)
       =(am,\eta(f\otimes m)+bn).
\]
A morphism $(\alpha,\beta)$ satisfies
$\beta\eta=\eta'(F\otimes\alpha)$.

\begin{proposition}[Triangular conversion]\label{conv:proposition}
Let $A$ be a finite-dimensional symmetric algebra over a field $\kk$
of characteristic two, let $F$ be a finite $A$-bimodule projective on
both sides, and let $X$ be a finite left $A$-module.  Suppose
$v\in\stHom{A}{X}{FX}$.  For each integer $a$, put
\[
 H^a=\stHom{A}{X}{X[a]},\qquad
 K^a=\stHom{A}{X}{(FX)[a]},
\]
and define
\begin{equation}\label{conv:delta}
 \delta^a:H^a\oplus H^a\longrightarrow K^a,
 \qquad (g,j)\longmapsto F(g)v+v[a]j.
\end{equation}
Assume that $\delta^0$ is surjective with nonzero kernel and that
$\delta^a$ is an isomorphism for every integer $a>0$.

Choose an actual map $v_0:X\to FX$ representing $v$ and an embedding
$i:X\to Q$ into a finite projective $A$-module.  Define
\begin{equation}\label{conv:finite-module}
 \Lambda=\begin{pmatrix}A&0\\F&A\end{pmatrix},\qquad
 Y=(FX\oplus Q)/\{(v_0(x),i(x)):x\in X\},\qquad
 Z=(X,Y,\iota),
\end{equation}
where $\iota:FX\to Y$ is induced by the first summand and specifies
the structure map of the left triangular module.  Then $Z$ is a finite
nonprojective Gorenstein-projective left $\Lambda$-module and
\[
 \Ext^a_\Lambda(Z,Z)=0=\Ext^a_\Lambda(Z,\Lambda)
 \qquad\text{for every integer }a>0.
\]
\end{proposition}

\subsection{Comparison of complete resolutions}

For homological complexes $P_\bullet,Q_\bullet$, use the full Hom complex
\begin{equation}\label{conv:full-hom}
 \Hom_R^a(P_\bullet,Q_\bullet)
   =\prod_{j\in\mathbb Z}\Hom_R(P_j,Q_{j-a}),
 \qquad
 (\partial f)_j=d^Q_{j-a}f_j+f_{j-1}d^P_j.
\end{equation}
This product convention is also used in \cite[Section 1.1.1]{Veliche2006}.
The product permits comparison maps with infinitely many nonzero
components.  Each component of the differential uses only two maps.
Its degree-$a$ cocycles are chain maps $P_\bullet\to Q_\bullet[a]$,
and its boundaries are nullhomotopic chain maps.  Throughout this section the
field has characteristic two, so the displayed signs apply to shifts
and cones as well.

\begin{lemma}[Complete-resolution comparison]\label{conv:comparison}
Let $R$ be a finite-dimensional algebra over a field of characteristic
two.  Let $P_\bullet$ be a totally acyclic complex of finite projective
left $R$-modules, and let $Q_\bullet$ be an exact complex of finite
projective left $R$-modules.  Put
\[
 C_j(P_\bullet)=\coker(d^P_{j+1}:P_{j+1}\to P_j),\qquad
 M=C_0(P_\bullet),\qquad N=C_0(Q_\bullet).
\]
There are natural, composition-compatible identifications
\[
 H^0\Hom_R^\bullet(P_\bullet,Q_\bullet)
       \simeq\stHom{R}{M}{N}.
\]
Applying this identification to shifted complexes gives, for every
integer $a$,
\begin{equation}\label{conv:shift-comparison}
 H^a\Hom_R^\bullet(P_\bullet,Q_\bullet)
 \simeq\stHom{R}{C_a(P_\bullet)}{N}
 \simeq\stHom{R}{M}{C_{-a}(Q_\bullet)}.
\end{equation}
For every integer $a>0$, these groups are naturally
$\Ext_R^a(M,N)$.
\end{lemma}

\begin{proof}
Exactness identifies $C_j(P_\bullet)$ with $\im d^P_j\subset P_{j-1}$.
We first record the extension property supplied by total acyclicity:
for every finite projective $T$, a map
$\phi:\im d^P_j\to T$ extends to $P_{j-1}$.  Indeed,
$\phi d^P_j:P_j\to T$ is a cocycle in $\Hom_R(P_\bullet,T)$.
Its exactness gives $\psi:P_{j-1}\to T$ with
$\psi d^P_j=\phi d^P_j$, which is the required extension.
This argument does not assume that projective $R$-modules are injective.

A chain map induces a map of cokernels.  Conversely, given $u:M\to N$,
first lift $u\epsilon_P$ to $f_0:P_0\to Q_0$, where $\epsilon_P$ and
$\epsilon_Q$ denote the quotient maps.  Projectivity and exactness give
$f_1,f_2,\ldots$ successively, with
$d^Q_jf_j=f_{j-1}d^P_j$.  To extend in the other direction,
$d^Q_0f_0$ kills $\im d^P_1=\ker d^P_0$.  It therefore induces a map
$\im d^P_0\to Q_{-1}$, which extends to $f_{-1}:P_{-1}\to Q_{-1}$ by
the preceding extension property.  The same argument gives
$f_{-2},f_{-3},\ldots$.  This constructs a full comparison map.

A nullhomotopy has components $h_j:P_j\to Q_{j+1}$ and satisfies
\[
 f_j=d^Q_{j+1}h_j+h_{j-1}d^P_j.
\]
On cokernels its first summand vanishes and its second factors through
$P_{-1}$.  Thus a nullhomotopic map induces a projective factorization.
For the converse, suppose the induced map factors as $M\to T\to N$
with $T$ finite projective.  Extend $M\to T$ to $P_{-1}$ by the extension
property, and lift $T\to N$ to $Q_0$.  Their composite is a homotopy
component $h_{-1}$ whose boundary induces exactly the given
factorization.  Subtract this boundary.  It remains to nullhomotope a
chain map $f'$ inducing zero on cokernels.

For this remaining map take $h_{-1}=0$.  Since $f'_0$ has image in
$\im d^Q_1$, projectivity gives $h_0:P_0\to Q_1$ with
$d^Q_1h_0=f'_0$.  Proceeding upwards, the residual
$f'_j+h_{j-1}d^P_j$ lands in $\ker d^Q_j=\im d^Q_{j+1}$ and hence
lifts to $h_j$.  Proceeding downwards from $j=-1$, the residual
$f'_j+d^Q_{j+1}h_j$ kills $\im d^P_{j+1}$ by the homotopy equation
already obtained at $j+1$.  It descends to $\im d^P_j$ and extends to
$h_{j-1}:P_{j-1}\to Q_j$ by the extension property.  Each recursion
assigns one component at a time and changes only adjacent equations.
Together they define a homotopy in the full product
\eqref{conv:full-hom}.  This proves the degree-zero identification.
Passing to cokernels respects composition, and the proved uniqueness
modulo homotopy gives its asserted compatibility.

Apply the degree-zero result to $P_\bullet[-a]$ or to
$Q_\bullet[a]$ to obtain \eqref{conv:shift-comparison}.  For $a>0$,
$C_a(P_\bullet)$ is the $a$th syzygy of $M$.  The projective-resolution
definition of $\Ext_R^a(M,N)$ quotients
$\Hom_R(C_a(P_\bullet),N)$ by maps extending to $P_{a-1}$.
These maps factor through a projective.  Conversely, every projective
factorization extends to $P_{a-1}$ by the extension property.  The two
quotients coincide, including when $a=1$.
\end{proof}

\subsection{Proof of the conversion principle}


\begin{proof}[Proof of Proposition~\ref{conv:proposition}]
We first check finiteness and introduce the two column functors.
Since $i$ is injective, the
map $\iota$ in \eqref{conv:finite-module} is injective and there is an
exact sequence
\begin{equation}\label{conv:finite-extension}
 0\longrightarrow FX\xrightarrow{\iota}Y
   \longrightarrow Q/i(X)\longrightarrow0.
\end{equation}
In particular all the modules in \eqref{conv:finite-module} are finite.
The algebra is finite as well, with
$\dim_\kk\Lambda=2\dim_\kk A+\dim_\kk F$.

The column functors below are standard; see \cite[Section 2]{EshraghiEtAl2014}.
Define
\[
 E(M)=(M,FM,\id),\qquad J(N)=(0,N,0).
\]
The functor $E$ is exact because $F$ is projective as a right
$A$-module; $J$ is exact directly.  Both preserve finite projectives:
$E(A)$ and $J(A)$ are the two projective columns of $\Lambda$, and
$\Lambda=E(A)\oplus J(A)$ as a left module.  The morphism condition
above gives
\begin{equation}\label{conv:four-homs}
\begin{aligned}
 \Hom_\Lambda(E(M),E(M'))&\simeq\Hom_A(M,M'),\\
 \Hom_\Lambda(J(N),J(N'))&\simeq\Hom_A(N,N'),\\
 \Hom_\Lambda(E(M),J(N))&=0,\\
 \Hom_\Lambda(J(N),E(M))&\simeq\Hom_A(N,FM).
\end{aligned}
\end{equation}
In the first line the lower component of a map is $F(\alpha)$.
The third line is a vanishing of ordinary module maps, so it holds
termwise for complexes in every shift.

Choose a complete resolution $P_\bullet$ of $X$ with
$P_{-1}=Q$ and $d_0=i\epsilon_P$.  This is possible by splicing a
projective resolution with a projective-injective coresolution, since
$A$ is symmetric.  Put
\[
 G_0=E(P_\bullet),\qquad G_1=J(P_\bullet).
\]
They are exact complexes of finite projective $\Lambda$-modules.
Moreover \eqref{conv:four-homs} gives
\begin{align*}
 \Hom_\Lambda(G_0,\Lambda)&\simeq\Hom_A(P_\bullet,A),\\
 \Hom_\Lambda(G_1,\Lambda)&\simeq
       \Hom_A(P_\bullet,F)\oplus\Hom_A(P_\bullet,A).
\end{align*}
These complexes are exact: the left $A$-modules $F$ and $A$ are
projective-injective.  Thus both $G_0$ and $G_1$ are totally acyclic.
Also $F\otimes_A P_\bullet$ is a complete resolution of $FX$.
Here right projectivity of $F$ gives exactness, left projectivity gives
projective terms, and self-injectivity of $A$ gives Hom exactness.

We can now compute all maps between the two induced complexes.
Lemma~\ref{conv:comparison} and \eqref{conv:four-homs} identify their
full Hom cohomology as
\begin{equation}\label{conv:column-cohomology}
\begin{aligned}
 H^a\End^\bullet_\Lambda(G_0)&\simeq H^a,
 &H^a\End^\bullet_\Lambda(G_1)&\simeq H^a,\\
 H^a\Hom^\bullet_\Lambda(G_1,G_0)&\simeq K^a,
 &\Hom^\bullet_\Lambda(G_0,G_1)&=0.
\end{aligned}
\end{equation}
Lift the actual map $v_0$ to a chain map $V:G_1\to G_0$ using the
comparison construction, and form
\[
 P_Z=\operatorname{Cone}(V),\qquad
 (P_Z)_j=(G_0)_j\oplus(G_1)_{j-1},\qquad
 d_j^{P_Z}=\begin{pmatrix}d^{G_0}&V_{j-1}\\0&d^{G_1}\end{pmatrix}.
\]
The degreewise split cone sequence shows that $P_Z$ is totally acyclic
with finite projective terms.

Its cokernel at degree zero is precisely the module $Z$ already
defined in \eqref{conv:finite-module}.  Its upper component is $X$.
For the lower component, first quotient $F\otimes_A P_0$ by the
image of $F\otimes_A d_1$, obtaining $FX$.  The remaining relations,
from $P_0$, are $(v_0(x),i(x))$: the lift $V_0$ induces $v_0$ on
cokernels and $d_0=i\epsilon_P$.  The resulting lower component is
therefore $Y$, with structure map $\iota$.  This supplies a complete
resolution of a concrete finite module.

It remains to compute its self-extensions.  Every graded endomorphism
of $P_Z$ has two diagonal blocks and one block from $G_1[1]$ to $G_0$;
the other direction vanishes by \eqref{conv:four-homs}.  Projection
to the diagonal gives a short exact sequence of cochain complexes
\begin{equation}\label{conv:end-sequence}
\begin{split}
0\longrightarrow\Hom^\bullet_\Lambda(G_1[1],G_0)
 \longrightarrow\End^\bullet_\Lambda(P_Z)
 \longrightarrow
 \End^\bullet_\Lambda(G_0)\oplus
 \End^\bullet_\Lambda(G_1[1])\longrightarrow0.
\end{split}
\end{equation}
This sequence is split in each graded degree, although the splitting
need not commute with the differentials.  Reindexing the full product
in \eqref{conv:full-hom} gives
\[
 H^a\Hom^\bullet_\Lambda(G_1[1],G_0)=K^{a-1}.
\]
The connecting map from diagonal degree $a$ to off-diagonal degree
$a+1$ is exactly \eqref{conv:delta}.  Indeed the differential of a
diagonal lift has cross term obtained by composing with $V$ on the
two sides.  The $E$-diagonal acts there by $F(g)$, and the $J$-diagonal
acts by precomposition with $j$, as specified in
\eqref{conv:four-homs}.  Compatibility of comparison with composition
identifies this cross term with $F(g)v+v[a]j$.

The long exact sequence of \eqref{conv:end-sequence} consequently
contains, for every integer $a$, the short exact sequence
\begin{equation}\label{conv:end-cohomology}
 0\longrightarrow\coker\delta^{a-1}
 \longrightarrow H^a\End^\bullet_\Lambda(P_Z)
 \longrightarrow\ker\delta^a\longrightarrow0.
\end{equation}
For $a=1$, surjectivity of $\delta^0$ and injectivity of $\delta^1$
make both outer terms zero.  For $a>1$, the two comparison maps are
isomorphisms.  Therefore the middle group vanishes for every $a>0$.
Lemma~\ref{conv:comparison} identifies it with
$\Ext^a_\Lambda(Z,Z)$.

At degree zero, \eqref{conv:end-cohomology} shows that
$H^0\End^\bullet_\Lambda(P_Z)$ surjects onto the nonzero kernel of
$\delta^0$.  By comparison this is a nonzero stable endomorphism
group of $Z$.  Its stable identity is consequently nonzero, since a
zero identity would annihilate every stable endomorphism.  Thus $Z$
is not projective.

Finally, total acyclicity gives exactness of
$\Hom_\Lambda(P_Z,\Lambda)$.  Its positive degrees compute the
ordinary groups $\Ext^a_\Lambda(Z,\Lambda)$ on the nonnegative
resolution tail, so these vanish for every $a>0$.  Exactness also
identifies $Z=\coker d_1^{P_Z}$ with the cycle
$\im d_0^{P_Z}=\ker d_{-1}^{P_Z}$ in this totally acyclic complex.
Hence $Z$ is Gorenstein-projective.
\end{proof}

Proposition~\ref{conv:proposition} reduces the construction to a finite
bimodule $F$ projective on both sides and a stable map $v$ with the stated two
action properties.  Surjectivity at degree zero is needed separately:
it is the term that removes a possible self-extension in degree one
in \eqref{conv:end-cohomology}.
